% Compile with pdfLaTeX (no external files or bibliography processor needed).
% Keep the submitted proposal to ONE page, including references.
\documentclass[11pt,letterpaper]{article}
\usepackage[T1]{fontenc}
\usepackage[utf8]{inputenc}
\usepackage[margin=0.75in]{geometry}
\usepackage{newtxtext,newtxmath}
\usepackage{microtype}
\usepackage{xcolor}
\usepackage{titlesec}
\usepackage{hyperref}

\definecolor{ProposalBlue}{HTML}{40619F}
\definecolor{SecondaryText}{HTML}{555555}
\hypersetup{colorlinks=true,urlcolor=ProposalBlue,linkcolor=ProposalBlue,
  pdfauthor={},pdftitle={New-Idea Proposal for a Paper Presentation}}
\pagestyle{empty}
\setlength{\parindent}{0pt}
\setlength{\parskip}{3pt}
\setlength{\emergencystretch}{1em}
\titleformat{\section}{\normalfont\bfseries\color{ProposalBlue}}{}{0pt}{}
\titlespacing*{\section}{0pt}{7pt}{2pt}
\newcommand{\proposalheader}[4]{%
  {\small\sffamily\bfseries\color{ProposalBlue}
   PAPER PRESENTATION\enspace\textbullet\enspace NEW-IDEA PROPOSAL}\par
  \vspace{3pt}
  {\fontsize{16}{18}\selectfont\bfseries #1\par}
  \vspace{4pt}
  \textbf{Student(s):} #2\hfill\textbf{Presentation:} #3\par
  \textbf{Based on:} #4\par
  \vspace{3pt}
  {\color{ProposalBlue}\hrule height 0.6pt}
  \vspace{4pt}
}
\newcommand{\guide}[1]{{\color{SecondaryText}#1}}
\newcommand{\wordbudget}[1]{{\normalfont\small\color{SecondaryText}\hfill #1 words}}
\newcommand{\referencesheading}{%
  \par\vspace{4pt}
  {\color{ProposalBlue}\hrule height 0.4pt}
  \vspace{3pt}
}

\begin{document}
\proposalheader{Question-Guided Worker Activation for MPLMs}
  {[Your name(s)]}{[Date]}
  {Liu et al., \textit{Message Passing Enables Efficient Reasoning} (2026) [1]}

\guide{\textit{Illustrative research extension; the experiments below are proposed, not completed.}}

\section*{1. Background and motivation}
Liu et al. introduce Message Passing Language Models (MPLMs), in which persistent
LLM threads exchange messages. In their long-context question-answering setup,
the parent launches workers for all document chunks, receives summaries, and then
requests evidence from selected workers [1, Sec.~5.3]. This motivates a cost question:
must every chunk be processed by an LLM before targeted reasoning begins?

\section*{2. Gap and research question}
Selective follow-up happens after the initial workers have processed their chunks.
When evidence lies in only a few chunks, this initial pass may be avoidable. I ask
whether choosing which workers to activate, rather than only whom to query later,
can reduce inference cost without missing evidence spread across documents.

\section*{3. Proposed idea and novelty}
I propose to rank chunks by keyword overlap with the question and initially activate
only the top two. These workers retain their contexts and use the existing message
protocol. If the parent identifies an unanswered subquestion, it reranks unopened
chunks using that subquestion and activates two more. Previously opened workers
remain available for follow-up. Expansion continues until the parent answers or a
shared generation budget is exhausted; unresolved questions can eventually trigger
full coverage. The added component is \emph{question-guided worker creation before LLM
processing}, not selective messaging itself. For example, a date lookup could start
with two relevant chunks instead of summarizing twenty.

\section*{4. Evaluation plan}
I will use Qwen3-30B-A3B, with 20 LongBench-v2 questions for development and 100
disjoint pilot-test questions sampled across task types and context lengths.
Compare original all-worker MPLM, fixed top-two activation, the proposed method,
and an ablation that opens chunks in random order. Hold model, chunking, decoding,
generation budget, maximum concurrency, and cache policy fixed. Measure answer
accuracy, mean wall-clock latency, processed input tokens, generated tokens, and
workers opened, including routing and retries. Run three seeds and report paired
confidence intervals. The pilot target is at least 20\% lower latency than all-worker
MPLM with accuracy within five percentage points; inconclusive intervals will be
reported as such.

\section*{5. Expected outcomes and risks}
I expect the largest savings when evidence is localized. Paraphrases and multi-hop
questions may defeat keyword ranking, while sequential expansion may erase the
benefit of fewer workers. Broader expansion is a fallback, not a correctness
guarantee. A negative result would help identify when processing every chunk at
the start is worth its cost.

\section*{6. Feasibility and deliverables}
Week 1: reproduce the baseline, check related work on retrieval-based routing,
and implement activation. Week 2: run the pilot and analyze missed-evidence errors.
Use a shared inference server without fine-tuning. Deliver a controller patch,
evaluation script, latency--accuracy plot, and representative failure cases.

\vfill
\referencesheading
{\small\textbf{Reference.} [1] X. Liu, D. Arora, G. Swamy, and A. Zanette.
\textit{Message Passing Enables Efficient Reasoning}. 2026.
\href{https://openreview.net/pdf?id=X5TrciQRnV}{OpenReview};
\href{https://arxiv.org/pdf/2607.01077}{arXiv full text}, Secs.~3 and 5.3.\par}
\end{document}
